\section{Aplicación de los Algoritmos al Problema de Portfolio}
\label{sec:aplicacion}

En este apartado se recogen los elementos comunes a todos los algoritmos empleados en la práctica. La idea es fijar primero cómo se representa una solución, cómo se evalúa y cómo se mantiene la factibilidad; después se describen los operadores compartidos de selección, cruce y mutación.

\subsection{Esquema de representación}
Se utiliza una \textbf{representación real} basada en un vector de pesos
$$
S = (s_1, s_2, \dots, s_N)
$$
donde cada componente $s_i$ indica la fracción del capital invertida en el activo $i$. Una solución es factible si se cumplen simultáneamente estas condiciones:
\begin{itemize}
    \item $\sum_{i=1}^N s_i = 1$.
    \item Para cada activo, o bien $s_i = 0$, o bien $s_i \in [lo, hi]$.
\end{itemize}

La población inicial se genera aleatoriamente y, a continuación, se aplica reparación para garantizar que todos los individuos parten de soluciones válidas.

\subsection{Función objetivo}
La calidad de una cartera se mide con una función de \textit{fitness} que combina beneficio y riesgo. Como en el problema se busca maximizar, una mayor aptitud indica una mejor solución\footnote{El cambio que se ha introducido es que se ha realizado la raíz para el cálculo del riesgo.}.

\begin{algorithm}[H]
\SetAlgoLined
\KwIn{Solución $S$, beneficios logarítmicos $logBenefits$, matriz de covarianza $Cov$, parámetro $\lambda$}
\KwOut{Valor de \textit{fitness} de $S$}

$beneficio \leftarrow 0$

\For{$i \leftarrow 1$ \KwTo $N$}{
    $beneficio \leftarrow beneficio + s_i \cdot logBenefits_i$
}

$riesgo \leftarrow 0$

\For{$i \leftarrow 1$ \KwTo $N$}{
    \For{$j \leftarrow 1$ \KwTo $N$}{
        $riesgo \leftarrow riesgo + s_i \cdot s_j \cdot Cov(i,j)$
    }
}

\Return $beneficio - \lambda \cdot \sqrt{\max(0, riesgo)}$
\caption{Evaluación de una solución}
\end{algorithm}
\textit{Nota: Se incorpora el operador $\max(0, \cdot)$ dentro de la raíz cuadrada para garantizar la estabilidad numérica del algoritmo, evitando valores negativos infinitesimales producidos por errores de precisión en punto flotante (IEEE 754).}

\subsection{Operador de reparación}
Los operadores genéticos pueden producir soluciones no factibles, por lo que después de cada cruce o mutación se aplica un operador de reparación. Su misión es devolver el vector al espacio de búsqueda permitido sin alterar innecesariamente la estructura de la solución\footnote{Se usa el mismo pseudocódigo de la práctica 1 ya que no ha sufrido cambios.}.

\newpage
\begin{algorithm}[H]
\SetAlgoLined
\KwIn{Vector de pesos inválido $w$, límites $lo$ y $hi$, tolerancia $\epsilon$}
\KwOut{Vector de pesos $w$ factible}
 \tcp{Fase 1: Normalización inicial}
 $suma \leftarrow \sum_{i=1}^{N} w_i$\;
 \If{$suma > 0$ y $|suma - 1.0| > \epsilon$}{
    $w_i \leftarrow w_i / suma \quad \forall i$\;
 }
 \tcp{Fase 2: Poda estricta de límites}
 \For{$i \leftarrow 1$ \KwTo $N$}{
    \If{$w_i < lo$}{ $w_i \leftarrow 0.0$ }
    \ElseIf{$w_i > hi$}{ $w_i \leftarrow hi$ }
 }
 \tcp{Fase 3: Redistribución iterativa}
 $suma \leftarrow \sum_{i=1}^{N} w_i$\;
 \While{$|suma - 1.0| > \epsilon$}{
    $\delta \leftarrow 1.0 - suma$\;
    \eIf{$\delta > 0$}{
        \tcp{Falta capital: repartir entre activos sin saturar}
        $C \leftarrow \{i \mid 0 < w_i < hi\}$\;
        \eIf{$C \neq \emptyset$}{
            $incremento \leftarrow \delta / |C|$\;
            \For{$i \in C$}{
                $w_i \leftarrow w_i + \min(incremento, hi - w_i)$\;
            }
        }{
            \tcp{Retorno: no hay activos con margen}
            Seleccionar un índice $k$ aleatorio donde $w_k == 0$\;
            $w_k \leftarrow lo$\;
        }
    }{
        \tcp{Sobra capital: recortar de activos activos}
        $C \leftarrow \{i \mid w_i > lo\}$\;
        $decremento \leftarrow |\delta| / |C|$\;
        \For{$i \in C$}{
            $w_i \leftarrow w_i - \min(decremento, w_i - lo)$\;
        }
    }
    $suma \leftarrow \sum_{i=1}^{N} w_i$\;
 }
 \Return $w$\;
 \caption{Operador de Reparación \texttt{fix(w)}}
 \label{alg:fix}
\end{algorithm}


\subsection{Operador de Vecindad (Trayectorias)}

Tanto la Búsqueda Local (BL) como el Enfriamiento Simulado (ES) utilizan el mismo operador de vecindad basado en la transferencia de capital entre activos.

\begin{algorithm}[H]
\SetAlgoLined
\KwIn{Solución $S$, ratio de transferencia $r=0.4$}
\KwOut{Solución vecina $S'$}

Seleccionar aleatoriamente un par de activos distintos $(i, j)$ de la cartera\;
\If{$s_i > 0$}{
    $transferencia \leftarrow s_i \cdot r$\;
    $s'_i \leftarrow s_i - transferencia$\;
    $s'_j \leftarrow s_j + transferencia$\;
    Reparar solución para asegurar límites $[lo, hi]$\;
}
\Return $S'$\;
\caption{Operador de vecindad por transferencia}
\end{algorithm}

\subsection{Operador de Mutación Brusca (ILS)}

El algoritmo ILS requiere una perturbación más fuerte que el operador de vecindad para saltar a otras regiones del espacio de búsqueda.

\begin{algorithm}[H]
\SetAlgoLined
\KwIn{Solución $S$, tasa de mutación $m=0.2$}
\KwOut{Solución mutada $S'$}

$N_{mut} \leftarrow \max(2, |S| \cdot m)$ activos elegidos al azar\;
Extraer los pesos $w_{idx}$ de esos activos\;
Barajar aleatoriamente (\textit{shuffle}) el conjunto de pesos extraídos\;
Reasignar los pesos barajados a los mismos activos originales\;
\Return $S'$\;
\caption{Mutación por barajado (ILS)}
\end{algorithm}